\documentclass[10pt,twocolumn]{article}
\usepackage[a4paper, left=1.5cm, right=1.5cm, top=2cm, bottom=3cm]{geometry}
\usepackage[T1]{fontenc}
\usepackage[utf8]{inputenc}
\usepackage[italian]{babel}
\usepackage{amsmath}
\usepackage{titling}
\usepackage{caption}
\usepackage{graphicx}
\usepackage{float}
\usepackage{relsize}
\usepackage{amsmath}
\usepackage{sectsty}
\usepackage{ragged2e}
\usepackage{circuitikz}
\usepackage{booktabs}
\usepackage{enumitem}
\usepackage{tikz}
\usepackage{physics}
\usepackage{xcolor}
\usepackage[most]{tcolorbox}
\usepackage{tikz-3dplot}
\usepackage{tikz}
\usepackage{ragged2e}
\usepackage{siunitx}
% \usepackage{booktabs}
\usepackage[colorlinks=true, linkcolor=black]{hyperref}  %per rendere l'indice genrale "interattivo"
\usepackage{enumitem}  %distanza degli itemize
\setlist[itemize]{itemsep=4pt, parsep=1pt}
\newtcolorbox{nota}{
  blanker,
  before skip=1em,
  after skip=1em,
  left=1em,
  borderline west={1pt}{0pt}{black},
  fontupper=\itshape,
  before upper={\noindent\textbf{Nota}:\quad}
}



\begin{document}
\justifying
	\title{\textbf{Verifica della legge dell'irradianza per una sorgente luminosa puntiforme}}
	\author{Brusini Alessio \hspace{0.7cm} Tedeschini Michele \hspace{0.7cm} Mirolo Manuele \hspace{0.7cm} Stroili Emanuele}
	\date{09 Ottobre 2026}
	\maketitle
	\newgeometry{left=3cm, right=3cm, top=4cm, bottom=4cm}
	\onecolumn
	\tableofcontents
\vspace{3cm}
	\begin{abstract}
		\centering
		\large
    L'esperimento consiste nell'ottenere la curva della tensione, misurata da un voltmetro a sua volta collegato 
    ad un fotodiodo, in funzione della distanza del fotodiodo rispetto alla sorgente: una lampadina. 
       
	\end{abstract}

	\newpage
\restoregeometry
\twocolumn


\section{legge teorica da verificare}
Data una sorgente di onde elettromagnetiche e date due superfici $S_{1}$ e $S_{2}$ poste ad un certa distanza l'una dall'altra e applicando
la legge di conservazione dell'energia si ottine la seguente relazione:
\begin{equation}
    \frac{I_1}{I_2} = \frac{r_1^2}{r_2^2} \quad \longrightarrow \quad I(r) = I_1 \left(\frac{r_1}{r}\right)^2
\end{equation}
Dove $r_{1}$,$r$ sono le distanze delle superici $S_{1}$ e $S_{2}$ dalla sorgenete.


\section{Apparato sperimentale}
\section{Grafici}
\section{Conclusione}
\end{document}